%!TEX root = ../../main.tex

\chapter{Ergebnisse und Auswertung}
\label{ch:ergebnisse}

Die in Kapitel~\ref{ch:durchfuehrung} beschriebene Durchführung hat eine geschlossene Datenbasis erzeugt, deren Auswertung Gegenstand dieses Kapitels ist. Es führt die Ergebnisse entlang der Struktur des Untersuchungsdesigns aus Kapitel~\ref{ch:konzeption} zusammen und bereitet damit die Beantwortung der beiden Forschungsfragen vor. Zunächst werden die Ergebnisse je \ac{ETL}-Aufgabe und Aufgabenkategorie dargestellt (Abschnitt~\ref{sec:ergebnisse-aufgaben}), anschließend die Modelle (Abschnitt~\ref{sec:vergleich-modelle}) und die Prompting-Strategien (Abschnitt~\ref{sec:vergleich-prompting}) gegeneinander verglichen. Darauf folgen die Ergebnisse der verketteten Pipeline mit der Trennung von eigener Schrittleistung und Fehlerfortpflanzung (Abschnitt~\ref{sec:ergebnisse-pipeline}) sowie die Auswertung der querschnittlichen Reproduzierbarkeitsdimension (Abschnitt~\ref{sec:ergebnisse-reproduzierbarkeit}). Abschnitt~\ref{sec:bestaetigungslauf} prüft sämtliche dieser Befunde anschließend auf einem zweiten, während der Entwicklung unbekannten Datensatz nach und beziffert dort zugleich Kosten, Token-Verbrauch und Laufzeit beider Ausführungsmodi. Abschnitt~\ref{sec:vergleich-baseline} stellt die \ac{KI}-gestützten Verfahren der regelbasierten Pipeline gegenüber, bevor Abschnitt~\ref{sec:diskussion} die Befunde bündelt, an den Hypothesen H1 und H2 aus Abschnitt~\ref{sec:forschungsdesign} spiegelt und in den Stand der Forschung einordnet. Die daraus abzuleitenden Empfehlungen sind Gegenstand des folgenden Kapitels~\ref{ch:handlungsempfehlungen}.

Sämtliche im Folgenden berichteten Genauigkeitswerte beziehen sich auf die in Abschnitt~\ref{subsec:genauigkeit} definierte Genauigkeit in ihrer abgeglichenen Ausprägung: Stimmt das Zielschema, wird auch bei abweichender Datensatzanzahl der Anteil korrekt reproduzierter Werte gemessen, statt das Ergebnis pauschal als unbrauchbar zu verwerfen. Fehlgeschlagene Ausführungen und strukturell nicht vergleichbare Ausgaben gehen mit dem Wert null in die Mittelwerte ein, da sie im praktischen Einsatz ebenfalls kein verwertbares Ergebnis liefern. Ergänzend wird als Erfolgsquote der Anteil derjenigen Läufe ausgewiesen, die überhaupt eine strukturell verwertbare Ausgabe erzeugt haben; die Differenz beider Kennzahlen trennt das Scheitern an der Ausführung von der inhaltlichen Abweichung eines lauffähigen Ergebnisses.

\section{Ergebnisse je ETL-Aufgabe}
\label{sec:ergebnisse-aufgaben}

Über alle 432 Läufe des Code-Generierungs-Modus hinweg erreichen die vier Modelle eine mittlere Genauigkeit von 0,573 bei einer Erfolgsquote von 62,5\,\%. Von den 432 Läufen führten 270 zu einer verwertbaren Ausgabe, 134 brachen mit einem Programmfehler ab und 28 lieferten eine lauffähige, aber strukturell abweichende Ausgabe. Der gesamte Aufwand belief sich auf 6,57~US-Dollar bei rund 194.000 Eingabe- und 583.000 Ausgabe-Token sowie einer mittleren Antwortzeit von 11,2~Sekunden je Aufruf. Diese Mittelwerte verdecken jedoch eine erhebliche Spreizung, die erst in der Aufschlüsselung nach Aufgaben sichtbar wird und die zentrale Beobachtung dieser Arbeit bereits vorwegnimmt: Die Ergebnisqualität variiert zwischen den Aufgaben deutlich stärker als zwischen den Modellen.

Abbildung~\ref{fig:uebersicht-heatmap} zeigt diese Spreizung im Gesamtüberblick. Aufgetragen sind alle neun Aufgaben gegen die vier Modelle in beiden Ausführungsmodi sowie die regelbasierte Vergleichspipeline, wobei für den Modus der Code-Generierung jeweils die beste der klassischen Prompting-Strategien ausgewiesen ist. Drei Muster treten unmittelbar hervor: Die formal beschreibbaren Aufgaben in der oberen Bildhälfte werden von allen Cloud-Modellen nahezu fehlerfrei gelöst; die semantisch geprägten Aufgaben zeigen ein uneinheitliches, zwischen den Modellen stark schwankendes Bild; und die zusammengesetzte Aufgabe aus Join, Bereinigung und Aggregation scheitert im Modus der Code-Generierung durchgängig, während sie im Direkt-Modus immerhin rund die Hälfte der Werte trifft.

\begin{figure}[htbp]
    \centering
    \includegraphics[width=\textwidth]{uebersicht_heatmap.png}
    \caption[Gesamtüberblick der Genauigkeit]{Abgeglichene Genauigkeit in Prozent über alle neun Aufgaben, vier Modelle und beide Ausführungsmodi sowie die regelbasierte Pipeline. Für die Code-Generierung ist jeweils die beste der drei klassischen Prompting-Strategien ausgewiesen.}
    \label{fig:uebersicht-heatmap}
\end{figure}

Tabelle~\ref{tab:ergebnisse-aufgaben} beziffert dieses Bild je Aufgabe. Ausgewiesen sind die über alle Modelle, Strategien und Wiederholungen gemittelte Genauigkeit, die zugehörige Erfolgsquote sowie zum Vergleich das Ergebnis der regelbasierten Pipeline. Die Reihenfolge folgt den drei Schwierigkeitsstufen aus Abschnitt~\ref{sec:etl-aufgaben}, sodass sich der Einfluss der Aufgabenkomplexität unmittelbar ablesen lässt.

\begin{table}[htbp]
    \centering
    \small
    \caption[Ergebnisse je ETL-Aufgabe]{Mittlere Genauigkeit und Erfolgsquote je Aufgabe über alle Modelle, Strategien und Wiederholungen des Code-Generierungs-Modus, gegenübergestellt der regelbasierten Pipeline.}
    \label{tab:ergebnisse-aufgaben}
    \begin{tabularx}{\textwidth}{@{}lXrrr@{}}
        \toprule
        \textbf{Kategorie} & \textbf{Aufgabe} & \textbf{Genauigkeit} & \textbf{Erfolg} & \textbf{Baseline} \\
        \midrule
        Bereinigung & Einfach -- Leerzeichen und fehlende Werte & 81,2\,\% & 81,2\,\% & 100,0\,\% \\
        Bereinigung & Mittel -- Datumsformate & 85,6\,\% & 89,6\,\% & 100,0\,\% \\
        Bereinigung & Schwer -- ISO-Ländercodes & 61,7\,\% & 66,7\,\% & 95,5\,\% \\
        \addlinespace
        Duplikate & Einfach -- exakte Duplikate & 77,1\,\% & 77,1\,\% & 100,0\,\% \\
        Duplikate & Mittel -- Schlüssel-Duplikate & 41,6\,\% & 41,7\,\% & 100,0\,\% \\
        Duplikate & Schwer -- unscharfe Duplikate & 19,1\,\% & 54,2\,\% & 96,9\,\% \\
        \addlinespace
        Transformation & Einfach -- Typkonvertierung & 68,0\,\% & 70,8\,\% & 100,0\,\% \\
        Transformation & Mittel -- abgeleitete Spalten & 81,2\,\% & 81,2\,\% & 100,0\,\% \\
        Transformation & Schwer -- Join und Aggregation & 0,0\,\% & 0,0\,\% & 12,1\,\% \\
        \bottomrule
    \end{tabularx}
\end{table}

Bemerkenswert ist zunächst, dass selbst die einfachsten Aufgaben im Mittel deutlich unter der Marke von 100~Prozent bleiben, obwohl jedes einzelne Modell sie mit gut gewähltem Prompt fehlerfrei löst. Diese Differenz geht überwiegend auf das lokal betriebene Modell und auf die schema-angereicherte Prompting-Variante zurück, deren Wirkung in den Abschnitten~\ref{sec:vergleich-modelle} und~\ref{sec:vergleich-prompting} gesondert betrachtet wird. Die Mittelwerte in Tabelle~\ref{tab:ergebnisse-aufgaben} sind daher nicht als Aussage über die erreichbare Qualität zu lesen, sondern als Aussage über die Verlässlichkeit einer beliebig herausgegriffenen Konfiguration: Sie messen, was zu erwarten ist, wenn Modell und Prompt nicht gezielt für die Aufgabe ausgewählt wurden.

\subsection{Datenbereinigung}
\label{subsec:ergebnisse-bereinigung}

Die Datenbereinigung ist die Kategorie mit der höchsten und zugleich stabilsten Ergebnisqualität. Das Entfernen überflüssiger Leerzeichen und die Kennzeichnung fehlender Werte lösen alle vier Modelle -- einschließlich des lokalen -- mit gut gewähltem Prompt vollständig korrekt, ebenso wie die Vereinheitlichung der vier Datumsformate, die bei allen Modellen zwischen 95,9 und 96,5~Prozent erreicht. Beide Aufgaben sind, wie in Abschnitt~\ref{sec:eingebaute-probleme} beschrieben, durch eine endliche Menge syntaktischer Regeln vollständig beschreibbar, und genau in diesem Bereich verhält sich der erzeugte Code so zuverlässig wie eine manuell geschriebene Lösung \autocite{coteDataCleaningMachine2024,rangineniReviewEnhancingData2023}.

Gerade hier zeigt sich jedoch auch ein Risiko, das sich in keiner Erfolgsquote niederschlägt. Skript~\ref{lst:pandas-object} stellt zwei Lösungen derselben Aufgabe gegenüber, die beide fehlerfrei durchlaufen und eine schemakonforme Ausgabe erzeugen. Die erste Variante wählt die zu bereinigenden Spalten über deren Datentyp aus -- ein in der Praxis verbreitetes und über Jahre korrektes Idiom. In der eingesetzten Pandas-Version werden Textspalten jedoch nicht mehr als generische Objektspalten, sondern als eigener Zeichenkettentyp geführt, sodass die Auswahl leer bleibt und die Bereinigung wirkungslos verpufft, ohne dass ein Fehler auftritt. Die zweite Variante benennt die betroffenen Spalten explizit und ist damit gegenüber dieser Änderung unempfindlich.

\begin{lstlisting}[language=Python, caption={Zwei lauffähige Lösungen derselben Bereinigungsaufgabe: Die obere Variante bleibt in der eingesetzten Pandas-Version wirkungslos, weil Textspalten nicht mehr als Objektspalten geführt werden; die untere benennt die Spalten explizit.}, label=lst:pandas-object]
# still wirkungslos: die Spaltenauswahl liefert keine Treffer
for col in df.select_dtypes(include=['object']).columns:
    df[col] = df[col].str.strip()

# wirksam: die betroffenen Spalten werden explizit benannt
for col in ["full_name", "email", "country", "registered_at"]:
    df[col] = df[col].astype(str).str.strip()
\end{lstlisting}

Dieser Befund ist für die Beurteilung der Eignung von erheblicher Bedeutung: Der Fehler ist nicht inhaltlicher, sondern umgebungsbedingter Natur, er erzeugt weder eine Fehlermeldung noch eine Schemaverletzung und wäre ohne eine bekannte Soll-Lösung im laufenden Betrieb nicht ohne Weiteres zu bemerken. Er entsteht dort, wo das Modell auf ein in seinen Trainingsdaten dominantes, inzwischen jedoch überholtes Idiom zurückgreift, und tritt damit systematisch und nicht zufällig auf \autocite{zhaoSurveyLargeLanguage2026,hassaniRoleChatGPTData2023}. Behoben wurde er im Rahmen der Prompt-Gestaltung nicht durch einen allgemeinen Hinweis auf Versionsunterschiede, sondern erst dadurch, dass dem Modell die Datentypen der Eingabespalten explizit mitgeteilt wurden.

Die semantische Vereinheitlichung der Länderangaben fällt mit 61,7~Prozent im Mittel deutlich ab, obwohl jedes Cloud-Modell mit gut gewähltem Prompt zwischen 94,6 und 95,8~Prozent erreicht. Der Grund für die verbleibende Lücke ist aufschlussreich und lässt sich anhand der Zwischenergebnisse der Pipeline exakt beziffern: Alle drei Cloud-Modelle lösen die Aufgabe, indem sie im erzeugten Code eine Zuordnungstabelle aufbauen und die Rohwerte darüber nachschlagen. Damit ist die Lösung nicht durch das Weltwissen des Modells begrenzt, sondern durch das, was es beim Schreiben des Codes vorab aufzählt. Nicht aufgelöst wurden auf diese Weise das Kürzel \texttt{SUI} für die Schweiz (14~Datensätze) und der Tippfehler \texttt{Deutshcland} (11~Datensätze) bei allen drei Modellen, ergänzt um \texttt{Großbritannien} bei einem und \texttt{OEsterreich} bei einem weiteren Modell (jeweils 7~Datensätze). Der Modus der Code-Generierung reduziert die semantische Aufgabe damit strukturell auf dieselbe Nachschlagelogik, die auch die regelbasierte Pipeline verwendet -- und erreicht folgerichtig ein ähnliches Ergebnis. Im Direkt-Modus, in dem das Modell jeden Wert einzeln beurteilt, statt eine verallgemeinernde Vorschrift zu formulieren, steigt die Genauigkeit auf 96,4 bis 96,5~Prozent und liegt damit erstmals über der regelbasierten Lösung \autocite{hassaniRoleChatGPTData2023,khanCognitiveETLUsing2025}.

\subsection{Duplikaterkennung}
\label{subsec:ergebnisse-duplikate}

Die Duplikaterkennung weist die stärkste Spreizung aller drei Kategorien auf, die Abbildung~\ref{fig:kategorie-dedup} nach Aufgaben und Strategien aufschlüsselt. Exakte Duplikate werden von allen vier Modellen und in jeder Strategie fehlerfrei entfernt; die Aufgabe ist mit einem einzigen Funktionsaufruf lösbar und trennt die Ansätze nicht. Bereits die schlüsselbasierte Entdoppelung anhand der E-Mail-Adresse fällt jedoch auf 41,6~Prozent im Mittel, und die unscharfen Duplikate erreichen mit 19,1~Prozent den niedrigsten Wert aller lösbaren Aufgaben.

\begin{figure}[htbp]
    \centering
    \includegraphics[width=\textwidth]{kategorie_deduplication_strategien.png}
    \caption[Duplikaterkennung je Strategie]{Duplikaterkennung: Genauigkeit je Modell und Prompting-Strategie im Vergleich zur Direkt-Verarbeitung und zur regelbasierten Pipeline (gestrichelte Linie).}
    \label{fig:kategorie-dedup}
\end{figure}

Bei näherer Betrachtung sind die niedrigen Mittelwerte weniger auf ein Unvermögen der Modelle als auf die Zweiseitigkeit der Aufgabe zurückzuführen. Anders als bei der Bereinigung, bei der jede Zeile für sich betrachtet wird, verlangt die Entdoppelung eine Entscheidung über das Verhältnis von Datensätzen zueinander, bei der sowohl zu wenig als auch zu viel entfernt werden kann. Die Ergebnisse zeigen beide Fehlerrichtungen: Bei den unscharfen Duplikaten liegt der beste Wert je Modell zwischen 41,5 und 69,4~Prozent, wobei ein zu enger Ähnlichkeitsschwellenwert echte Zwillinge übersieht, während ein zu weiter Schwellenwert unbeteiligte Datensätze zusammenführt. In der verketteten Pipeline trat dieser zweite Fehlertyp in besonders klarer Form auf: Werden akademische Grade und Anreden nicht als solche erkannt, sondern als Namensbestandteile behandelt, entstehen Vergleichsschlüssel, unter denen bis zu sieben verschiedene Kunden zusammenfallen. Erst die explizite Anweisung, Titel nicht als Teil des Namens zu behandeln und Übereinstimmung sowohl im Namen als auch in der E-Mail-Adresse zu verlangen, löste dieses Problem auf.

Der deutlichste Befund dieser Kategorie ergibt sich aus dem Vergleich der beiden Ausführungsmodi. Während die beste Code-Generierung bei den unscharfen Duplikaten 69,4~Prozent erreicht, liegt der Direkt-Modus bei denselben Modellen zwischen 93,5 und 96,5~Prozent. Das Muster entspricht dem der Länder-Vereinheitlichung und hat dieselbe Ursache: Die Beurteilung, ob die beiden in Tabelle~\ref{tab:dirty-duplikate} gegenübergestellten Datensätze \texttt{Uta Lange / gkabus@example.org} und \texttt{Uta\textvisiblespace\textvisiblespace Lange / gka.bus@example.org} dieselbe Person bezeichnen, ist für ein Sprachmodell eine unmittelbar zugängliche Einzelfallentscheidung, aber eine schwer zu verallgemeinernde Regel. Sobald das Modell diese Entscheidung in Code fassen muss, ist es gezwungen, sie in einen numerischen Schwellenwert zu übersetzen, und verliert dabei genau die Fähigkeit, die es der regelbasierten Pipeline überlegen macht \autocite{coteDataCleaningMachine2024,hassaniRoleChatGPTData2023}.

\subsection{Datentransformation}
\label{subsec:ergebnisse-transformation}

Typkonvertierung und abgeleitete Spalten fügen sich in das Bild der formal beschreibbaren Aufgaben ein: Beide werden von den Cloud-Modellen mit gut gewähltem Prompt vollständig oder nahezu vollständig korrekt gelöst, wobei die Typkonvertierung mit 68,0~Prozent im Mittel vor allem deshalb abfällt, weil das lokale Modell hier lediglich 26,7~Prozent erreicht und die schema-angereicherte Strategie zu abstürzenden Konvertierungen verleitet.

Die zusammengesetzte Aufgabe aus Join, Bereinigung und Aggregation ist dagegen die einzige der neun Aufgaben, die im Modus der Code-Generierung von keinem Modell und in keiner Strategie gelöst wurde. Von den 48 Läufen brachen 33 mit einem Programmfehler ab, die verbleibenden 15 erzeugten eine lauffähige, aber inhaltlich abweichende Ausgabe. Deren Analyse zeigt zwei wiederkehrende Fehlerbilder. Im ersten liefert das Ergebnis 65 statt der erwarteten 60 Gruppen und aggregiert lediglich 485 der 2.000 Bestellungen, weil die Verknüpfung Bestellungen ohne passenden Kundendatensatz verwirft. Im zweiten werden zwar alle 2.000 Bestellungen korrekt aggregiert, die Gruppierung erfolgt jedoch über 42 verschiedene, lediglich in Großbuchstaben überführte Rohschreibweisen der Länderangabe statt über die zehn ISO-Codes, sodass etwa \texttt{DE}, \texttt{DEUTSCHLAND}, \texttt{GERMANY}, \texttt{GER} und \texttt{DEUTSHCLAND} als fünf getrennte Länder erscheinen.

Beide Fehlerbilder sind aufschlussreich, weil sie nicht die Aggregation selbst betreffen. Die Modelle beherrschen Verknüpfung und Gruppierung als solche; sie scheitern daran, die vorgelagerte semantische Bereinigung im selben Arbeitsschritt mitzudenken. Dass die Genauigkeit dabei exakt auf null fällt, ist der Zusammensetzung der Aufgabe geschuldet: Weil die bereinigte Länderangabe zugleich Gruppierungsschlüssel ist, wirkt sich jede unvollständige Bereinigung auf sämtliche aggregierten Werte aus, sodass kein einziger Kennzahlenwert korrekt bleibt. Die Kennzahl misst hier folglich nicht die Fähigkeit zur Aggregation, sondern die Fähigkeit zur Komposition mehrerer Verarbeitungsschritte. Dass es sich um genau diese Unterscheidung handelt, belegt Abschnitt~\ref{sec:ergebnisse-pipeline}: Wird derselbe Schritt als Glied einer Kette mit bereits bereinigten Eingabedaten gestellt, lösen ihn alle vier Modelle fehlerfrei.

\section{Vergleich der Modelle}
\label{sec:vergleich-modelle}

Abbildung~\ref{fig:gesamt-modelle} fasst den Modellvergleich über alle 432 Läufe zusammen. Unter den kommerziell betriebenen Modellen liegen die Ergebnisse eng beieinander: Das Modell von OpenAI erreicht mit einer Genauigkeit von 0,733 bei einer Erfolgsquote von 79,6~Prozent den höchsten Wert, gefolgt von Anthropic mit 0,661 und 73,1~Prozent sowie Google mit 0,635 und 71,3~Prozent. Die Abstände zwischen diesen drei Modellen sind mit rund zehn Prozentpunkten deutlich geringer als die zuvor beobachteten Unterschiede zwischen den Aufgaben, die von 0 bis über 95~Prozent reichen. Für die Frage der Eignung ist damit die Wahl der Aufgabe erheblich folgenreicher als die Wahl des Anbieters -- ein Befund, der die in der Literatur verbreitete Fokussierung auf Modellranglisten relativiert \autocite{2026AIIndex,ComparisonAIModels}.

\begin{figure}[htbp]
    \centering
    \includegraphics[width=\textwidth]{gesamt_modelle.png}
    \caption[Modellvergleich über alle Läufe]{Modellvergleich über alle 432 Läufe der Code-Generierung: Genauigkeit und Erfolgsquote, Ausgang der je 108 Läufe sowie Gesamtkosten und mittlere Antwortzeit.}
    \label{fig:gesamt-modelle}
\end{figure}

Deutlich abgesetzt ist das lokal betriebene Modell der Llama-Reihe mit einer Genauigkeit von 0,261 und einer Erfolgsquote von 25,9~Prozent. Das mittlere Bild in Abbildung~\ref{fig:gesamt-modelle} zeigt, worin dieser Abstand besteht: Von seinen 108 Läufen brachen 74 mit einem Programmfehler ab, während die Cloud-Modelle zwischen 11 und 25 solcher Abbrüche aufweisen. Das lokale Modell scheitert damit weit überwiegend nicht an der inhaltlichen Aufgabe, sondern daran, überhaupt lauffähigen Code zu erzeugen. Diese Trennung ist praktisch bedeutsam, weil sie einen qualitativen Unterschied markiert: Ein inhaltlich abweichendes Ergebnis lässt sich prüfen und korrigieren, ein nicht lauffähiges Skript hingegen erzwingt einen erneuten Aufruf und verschiebt das Problem lediglich \autocite{zhaoSurveyLargeLanguage2026,hassaniRoleChatGPTData2023}. Dass das Modell die Aufgaben durchaus versteht, belegt der in Abschnitt~\ref{sec:ergebnisse-pipeline} dargestellte isolierte Kettendurchlauf, in dem es bei bis zu drei zugestandenen Versuchen je Schritt alle neun Schritte durchläuft.

Diese Beobachtung erfordert eine Einordnung, die zugleich eine Entscheidung über den weiteren Aufbau dieses Kapitels begründet. Der gemessene Abstand ist keine Aussage über lokal betreibbare Modelle im Allgemeinen, sondern über dasjenige Modell, das sich auf der für diese Arbeit verfügbaren Hardware sinnvoll ausführen ließ. Die Größe eines lokal betriebenen Modells ist unmittelbar durch den verfügbaren Grafikspeicher begrenzt, und mit ihr die Fähigkeit, längere Anweisungen zu verarbeiten und syntaktisch fehlerfreien Code zu erzeugen. Die kommerziellen Modelle unterliegen dieser Beschränkung nicht, da die Ausführung beim Anbieter stattfindet. Ein Vergleich zwischen beiden misst somit nicht allein die Modellqualität, sondern zu einem erheblichen Teil die Ausstattung des ausführenden Rechners \autocite{zhaoSurveyLargeLanguage2026,2026AIIndex}.

Aus diesem Grund wird das lokale Modell in den folgenden Vergleichen dieses Kapitels -- dem Vergleich der Prompting-Strategien, der verketteten Pipeline und dem Bestätigungslauf -- nicht weiter geführt. Der Grund ist nicht sein schwächeres Abschneiden als solches, sondern dass seine Ausfälle die Mittelwerte in einer Weise verschieben, die sich nicht mehr der untersuchten Einflussgröße zurechnen lässt: Wie Abschnitt~\ref{sec:vergleich-prompting} zeigt, entsteht die dort gemessene Rangfolge der Prompting-Strategien nahezu vollständig aus dem Verhalten dieses einen Modells. Seine Ergebnisse aus der Faktormatrix bleiben in Abschnitt~\ref{sec:ergebnisse-aufgaben} und im vorliegenden Abschnitt vollständig ausgewiesen; ausgeschlossen wird es nur dort, wo es eine andere Fragestellung überlagern würde.

Für die Praxis folgt daraus ausdrücklich keine Empfehlung gegen den lokalen Betrieb. Er ist die einzige der untersuchten Optionen, bei der keine Daten das eigene Netz verlassen, und er verursacht keine laufenden Nutzungskosten -- gegenüber den 2,05 bis 2,33~US-Dollar je Modell für die Faktormatrix und rund 0,19~US-Dollar je vollständigem Pipeline-Durchlauf ein Vorteil, der bei dauerhaftem Betrieb erheblich wird. Da die beobachtete Schwäche überwiegend in der Erzeugung lauffähigen Codes und nicht im inhaltlichen Verständnis der Aufgaben liegt, ist zu erwarten, dass ein auf leistungsfähigerer Hardware ausführbares, größeres Modell diesen Abstand deutlich verringert. Für Vorhaben mit entsprechender Ausstattung ist eine eigene Erprobung daher lohnend, und die in Abschnitt~\ref{sec:empfehlungen-auswahl} genannten Bedingungen -- kleinteiliger Aufgabenzuschnitt und mehrere zugestandene Versuche je Schritt -- geben dafür den Rahmen vor \autocite{khanLeveragingArtificialIntelligence2025,2026AIIndex}.

Die Kostenseite relativiert die Qualitätsunterschiede zusätzlich. Über alle 108 Läufe eines Modells liegen die Nutzungskosten zwischen 2,05 und 2,33~US-Dollar und damit in derselben Größenordnung; das lokal ausgeführte Modell verursacht überhaupt keine Nutzungskosten, benötigt dafür jedoch eigene Hardware und lieferte, wie in Abschnitt~\ref{sec:reproduzierbarkeit} beschrieben, im Direkt-Modus wegen des begrenzten Kontextfensters in den meisten Fällen kein verwertbares Ergebnis. Die Antwortzeiten unterscheiden sich mit 8,9 bis 13,6~Sekunden je Aufruf ebenfalls nur moderat. Ein Modell allein aufgrund von Kosten oder Geschwindigkeit auszuwählen, ist auf dieser Grundlage nicht zu begründen; die Auswahl folgt vielmehr der Zuverlässigkeit bei der konkret vorliegenden Aufgabenart.

\section{Vergleich der Prompting-Strategien}
\label{sec:vergleich-prompting}

Der Vergleich der vier Prompting-Strategien in Abbildung~\ref{fig:gesamt-strategien} liefert zunächst ein für die Ausgangserwartung überraschendes Ergebnis. Über alle vier Modelle und Aufgaben hinweg erzielt nicht eine der aufwändigeren Strategien, sondern das einfache Zero-Shot-Prompting mit 0,655 die höchste Genauigkeit, gefolgt vom Chain-of-Thought-Prompting mit 0,641 und dem Few-Shot-Prompting mit 0,609. Die schema-angereicherte Variante fällt mit 0,386 deutlich ab. Der Aufwand wächst dabei in die entgegengesetzte Richtung: Das Few-Shot-Prompting benötigt mit im Mittel 834~Eingabe-Token annähernd das Dreifache des Zero-Shot-Prompts mit 290~Token.

Diese Rangfolge hält jedoch einer genaueren Prüfung nicht stand, und die Ursache ist das in Abschnitt~\ref{sec:vergleich-modelle} beschriebene lokale Modell. Beschränkt man dieselbe Auswertung auf die drei kommerziellen Modelle, verschwindet der Unterschied nahezu vollständig: Zero-Shot erreicht 72,6, Few-Shot 73,3 und Chain-of-Thought 73,4~Prozent -- eine Spanne von 0,8 statt 4,6~Prozentpunkten. Das lokale Modell allein trägt den Effekt, und zwar in ausgeprägter Form: Es fällt von 44,1~Prozent unter dem Zero-Shot-Prompt auf 24,0~Prozent unter dem Few-Shot-Prompt, während seine Erfolgsquote von 44,4 auf 25,9~Prozent sinkt. Der dreifach längere Prompt überfordert es, was zu den in Abschnitt~\ref{sec:vergleich-modelle} genannten Grenzen der verfügbaren Hardware passt. Die scheinbar allgemeine Aussage \enquote{der einfachste Prompt ist der beste} ist damit keine Aussage über Prompting-Strategien, sondern über die Kapazität eines bestimmten Modells \autocite{changEfficientPromptingMethods2024,schulhoffPromptReportSystematic2025}.

Aussagekräftig ist der Strategievergleich folglich erst auf Aufgabenebene und unter den kommerziellen Modellen. Dort zeigt sich, dass sechs der neun Aufgaben zwischen den Strategien überhaupt nicht unterscheiden: Die Bereinigungsaufgaben und die exakten Duplikate werden von jeder Strategie nahezu fehlerfrei gelöst, die zusammengesetzte Aufgabe von keiner. Der gesamte messbare Unterschied konzentriert sich auf die beiden anspruchsvollen Entdoppelungsaufgaben, und dort mit klarer Zuordnung: Bei den unscharfen Duplikaten führt das Few-Shot-Prompting mit 41,5 gegenüber 30,7~Prozent für Zero-Shot und 21,3 für Chain-of-Thought, bei den Schlüssel-Duplikaten hingegen das Chain-of-Thought-Prompting mit 51,1 gegenüber jeweils 33,3~Prozent. Beispiele wirken damit dort, wo eine Ähnlichkeitsentscheidung zu treffen ist, für die eine verbale Beschreibung schlecht taugt; die Aufforderung zur schrittweisen Herleitung wirkt dort, wo eine mehrstufige Auswahlregel umzusetzen ist. Wie Abschnitt~\ref{subsec:bestaetigung-strategien} zeigt, ist genau dieses aufgabenbezogene Muster auf dem unbekannten Datensatz reproduzierbar, die Gesamtrangfolge dagegen nicht \autocite{sahooSystematicSurveyPrompt,liuPretrainPromptPredict2023}.

\begin{figure}[htbp]
    \centering
    \includegraphics[width=\textwidth]{gesamt_strategien.png}
    \caption[Vergleich der Prompting-Strategien]{Prompting-Strategien im Vergleich: Genauigkeit und Erfolgsquote über alle Läufe, Genauigkeit je Aufgabenkategorie sowie mittlere Prompt-Länge.}
    \label{fig:gesamt-strategien}
\end{figure}

Das mittlere Bild der Abbildung, das die Kategorien über alle vier Modelle zusammenfasst, ist entsprechend mit Vorsicht zu lesen: Der dort bei der Datenbereinigung sichtbare Vorsprung des Zero-Shot-Promptings von 89,2 gegenüber 80,0 und 75,3~Prozent geht auf dieselbe Ursache zurück und besteht unter den kommerziellen Modellen nicht. Festzuhalten bleibt der allgemeinere Zusammenhang, der sich in beiden Betrachtungen zeigt: Der Nutzen zusätzlicher Anleitung ist nicht generell, sondern an die Eigenart der Aufgabe gebunden. Wo eine Aufgabe in einem einzigen gedanklichen Schritt lösbar ist, stiftet sie keinen Mehrwert und verleitet zu unnötig komplexem und damit fehleranfälligerem Code; wo mehrere aufeinander aufbauende Entscheidungen zu treffen sind oder eine Ermessensgrenze zu setzen ist, wirkt sie \autocite{sahooSystematicSurveyPrompt,schulhoffPromptReportSystematic2025}.

Der Einbruch der schema-angereicherten Strategie verdient gesonderte Betrachtung, da er dem eigentlich beabsichtigten Zweck zuwiderläuft. Die Angabe des Zielschemas mit Spaltennamen und Datentypen sollte die strukturelle Korrektheit der Ausgabe unterstützen; tatsächlich fällt die Genauigkeit bei der Duplikaterkennung auf 14,3~Prozent und damit auf ein Viertel des Werts der übrigen Strategien. Die Ursache liegt darin, dass die Modelle die genannten Zieltypen als unmittelbaren Handlungsauftrag auffassen und Konvertierungen vornehmen, die die Aufgabenstellung gar nicht verlangt -- etwa das Umwandeln eines Datumsfeldes, dessen Rohwerte in dieser Aufgabe noch in vier verschiedenen Formaten vorliegen und die Konvertierung folglich zum Abbruch bringen. Zusätzlicher Kontext ist damit nicht per se hilfreich: Er wirkt nur, wenn er zur Aufgabe passt, und richtet Schaden an, wenn er über sie hinausweist \autocite{changEfficientPromptingMethods2024,liuPretrainPromptPredict2023}.

Zusammenfassend ist die Wirkung der Prompt-Gestaltung damit anders zu beurteilen, als es die Gesamtrangfolge nahelegt. Sie ist nicht gleichmäßig über die Aufgaben verteilt, sondern auf wenige konzentriert: Bei sechs der neun Aufgaben ist sie ohne messbare Wirkung, bei zwei entscheidet sie über bis zu 20~Prozentpunkte, und bei einer -- der schema-angereicherten Variante in der Duplikaterkennung -- richtet sie erheblichen Schaden an. Ihr Gewicht ist damit dort, wo sie überhaupt wirkt, durchaus mit dem der Modellwahl vergleichbar; eine allgemeine Empfehlung für oder gegen aufwändigere Prompts lässt sich daraus jedoch nicht ableiten, sondern nur eine je Aufgabenart \autocite{schulhoffPromptReportSystematic2025,changEfficientPromptingMethods2024}.

\section{Ergebnisse der verketteten Pipeline}
\label{sec:ergebnisse-pipeline}

Die bisher dargestellten Ergebnisse betrachten jede Aufgabe für sich. Die in Abschnitt~\ref{sec:forschungsdesign} eingeführte dritte Untersuchungsebene prüft, wie sich dieses Bild verändert, wenn dieselben Aufgaben als durchgängige Strecke ausgeführt werden, in der jeder Schritt auf dem Ergebnis seines Vorgängers aufsetzt.

\subsection{Qualitätsverlauf über die Kette}
\label{subsec:pipeline-verlauf}

Abbildung~\ref{fig:pipeline-verlauf} zeigt den Verlauf der neun Schritte für alle vier Modelle und die regelbasierte Pipeline. Drei der vier Modelle durchlaufen die Kette vollständig und erreichen in den Schritten~1 bis~8 durchweg Werte zwischen 97,6 und 100~Prozent. Der Aufwand einer vollständigen Strecke ist dabei gering: Er beträgt je Modell zwischen 0,176 und 0,198~US-Dollar sowie 117 bis 153~Sekunden, gegenüber 6,57~US-Dollar für die vollständige Faktormatrix. Die regelbasierte Pipeline benötigt für dieselbe Strecke 0,17~Sekunden und keine Nutzungskosten.

\begin{figure}[htbp]
    \centering
    \includegraphics[width=0.92\textwidth]{pipeline_verlauf.png}
    \caption[Qualitätsverlauf der verketteten Pipeline]{Abgeglichene Genauigkeit je Schritt der verketteten Pipeline. Grau markierte Felder kennzeichnen Schritte, die infolge eines vorangegangenen Abbruchs blockiert wurden.}
    \label{fig:pipeline-verlauf}
\end{figure}

Der Endschritt bricht dagegen bei allen Teilnehmern ein: Das Modell von Google erreicht 81,8~Prozent, jenes von OpenAI 72,7~Prozent und jenes von Anthropic 71,2~Prozent, die regelbasierte Pipeline 60,6~Prozent. Bemerkenswert ist zunächst, dass die Werte damit weit über dem Ergebnis derselben Aufgabe als Einzellauf liegen, wo sie bei null lag. Die Zerlegung des Endschritts erklärt beides: Wendet man die Referenzlogik der Aggregation auf die tatsächlich gelieferten Zwischenergebnisse jedes Teilnehmers an, ergibt sich für alle vier eine Übereinstimmung von 1,000. Der von jedem Teilnehmer selbst erzeugte Aggregationscode ist also nachweislich korrekt; der Verlust entsteht ausschließlich aus den Daten, die er verarbeitet. Konkret bleiben bei den Modellen 5 bis 8 und bei der regelbasierten Pipeline 20 Bestellungen ohne zugehörigen Kundendatensatz zurück, weil in der vorgelagerten Entdoppelung jeweils der andere von zwei unscharfen Zwillingen behalten wurde. Da die Aggregation über Land und Kategorie gruppiert, macht jede solche Bestellung zwei Gruppen in beiden Kennzahlen falsch, sodass wenige fehlplatzierte Datensätze einen überproportionalen Ausschlag erzeugen.

Abbildung~\ref{fig:pipeline-vs-einzellauf} stellt beide Betrachtungsweisen unmittelbar gegenüber. Für sechs der neun Schritte liegt das Ergebnis in der Kette über dem des unabhängigen Einzellaufs, teilweise erheblich: bei den unscharfen Duplikaten um 54,3, bei der Aggregation um 75,2 und bei den Schlüssel-Duplikaten um 35,3~Prozentpunkte. Dieser zunächst widersprüchlich wirkende Befund hat eine klare Ursache: In der Kette erhält jeder Schritt bereits vorbereitete Daten und eine eng umrissene Teilaufgabe, während er im Einzellauf zusätzlich alle vorgelagerten Mängel mitbehandeln muss. Die Zerlegung einer Strecke in kleine, klar abgegrenzte Schritte verbessert das Ergebnis somit unabhängig vom eingesetzten Modell -- ein Befund, der unmittelbar auf die praktische Gestaltung \ac{KI}-gestützter \ac{ETL}-Strecken durchschlägt \autocite{khanOverviewETLTechniques2024,nagabhyruBridgingTraditionalETL2022}.

\begin{figure}[htbp]
    \centering
    \includegraphics[width=\textwidth]{pipeline_vs_einzellauf.png}
    \caption[Verkettete Pipeline gegen unabhängige Einzelläufe]{Dieselben Aufgaben als Glied einer Kette und als unabhängiger Einzellauf, gemittelt über die Modelle. Die Kette schneidet in sechs von neun Schritten besser ab.}
    \label{fig:pipeline-vs-einzellauf}
\end{figure}

\subsection{Trennung von eigener Leistung und Fehlerfortpflanzung}
\label{subsec:pipeline-fortpflanzung}

Der bis hierhin gemessene Verlauf lässt offen, welcher Anteil eines schwachen Schrittergebnisses dem Schritt selbst und welcher den Vorgängerschritten zuzurechnen ist. Der in Abschnitt~\ref{sec:forschungsdesign} vorgesehene isolierte Durchlauf, in dem jeder Schritt die Soll-Ausgabe seines Vorgängers erhält, beantwortet diese Frage. Abbildung~\ref{fig:pipeline-fortpflanzung} stellt beide Durchläufe je Modell und Schritt gegenüber; die Differenz beider Balken beziffert unmittelbar den Preis der Fehlerfortpflanzung.

\begin{figure}[htbp]
    \centering
    \includegraphics[width=\textwidth]{pipeline_fortpflanzung.png}
    \caption[Fehlerfortpflanzung in der Pipeline]{Dieselbe Aufgabe mit fehlerfreier Eingabe (hell) und in der Kette (dunkel). Rote Kreuze kennzeichnen abgebrochene oder blockierte Schritte.}
    \label{fig:pipeline-fortpflanzung}
\end{figure}

Das Ergebnis ist eindeutig: Mit fehlerfreier Eingabe lösen alle vier Modelle den Endschritt vollständig korrekt. Der in der Kette gemessene Verlust von 18,2 bis 28,8~Prozentpunkten ist damit zu hundert Prozent geerbt und nicht dem Endschritt anzulasten. Dasselbe gilt in abgeschwächter Form für die Entdoppelungsschritte, die isoliert zwischen 98,8 und 100~Prozent erreichen und in der Kette ein bis zwei Prozentpunkte verlieren. Als eigene, nicht geerbte Schwäche verbleibt allein die semantische Länder-Vereinheitlichung: Sie ist der einzige Schritt, den auch mit fehlerfreier Eingabe kein Modell vollständig löst (99,1 / 98,8 / 98,8 / 82,3~Prozent), und zwar aus dem in Abschnitt~\ref{subsec:ergebnisse-bereinigung} belegten Grund der unvollständig aufgezählten Zuordnungstabelle.

Auch das Verhalten des lokalen Modells klärt sich erst durch diese Trennung auf. In der Kette erreicht es 6 von 9~Schritten und bricht in Schritt~5 ab; isoliert durchläuft es alle neun Schritte und löst den Endschritt ebenfalls fehlerfrei. Die Ursache des Abbruchs liegt drei Schritte früher: Seine Lösung der Datumsvereinheitlichung ließ 145 der 540~Registrierungsdaten unverändert als Unix-Zeitstempel stehen, was für sich genommen eine Genauigkeit von 94,6~Prozent bedeutet und die Kette nicht unterbricht. Erst der Entdoppelungsschritt, der auf diesem Feld eine Datumskonvertierung vornimmt, scheitert an den verbliebenen Zeitstempeln und blockiert die beiden folgenden Schritte. Zwei der drei Ausfälle des lokalen Modells sind damit nicht eigene Fehler, sondern Folgen eines früheren, für sich betrachtet unauffälligen Ergebnisses.

Für die Praxis folgt daraus eine unmittelbar umsetzbare Konsequenz: In verketteten Strecken entscheidet die Qualität der frühen, semantisch geprägten Schritte über das Endergebnis, und ein Fehler wird umso teurer, je weiter vorn er entsteht. Eine Prüfung nach jedem Schritt ist deshalb wirksamer als eine Endkontrolle, weil sie einen Mangel dort sichtbar macht, wo er noch mit begrenztem Aufwand zu beheben ist, statt ihn erst in seiner verstärkten Form am Ende der Strecke zu entdecken \autocite{ridzuanReviewDataQuality2024,boulahiaMulticriteriaEvaluationResearch2024}.

\section{Reproduzierbarkeit der Ergebnisse}
\label{sec:ergebnisse-reproduzierbarkeit}

Die in Abschnitt~\ref{sec:forschungsdesign} als querschnittliche Dimension verankerte Reproduzierbarkeit wird aus der Streuung der Genauigkeit über die drei Wiederholungen jeder Konfiguration bestimmt. Abbildung~\ref{fig:gesamt-reproduzierbarkeit} zeigt das Ergebnis. Über alle Konfigurationen hinweg beträgt die mittlere Standardabweichung 0,041, und 74,3~Prozent der Konfigurationen liefern über alle drei Wiederholungen hinweg ein exakt identisches Ergebnis. Die niedrig gewählte Temperatur hat die Variabilität somit wirksam begrenzt, aber nicht beseitigt: In gut einem Viertel der Fälle führt derselbe Prompt bei demselben Modell zu messbar unterschiedlichen Ergebnissen.

\begin{figure}[htbp]
    \centering
    \includegraphics[width=\textwidth]{gesamt_reproduzierbarkeit.png}
    \caption[Reproduzierbarkeit der Ergebnisse]{Streuung der Genauigkeit über die drei Wiederholungen je Konfiguration, aufgeschlüsselt nach Aufgabe und nach Modell.}
    \label{fig:gesamt-reproduzierbarkeit}
\end{figure}

Die Verteilung dieser Instabilität ist dabei nicht zufällig. Sie konzentriert sich auf genau jene Aufgaben, die auch inhaltlich am schwersten fallen: Die unscharfen Duplikate weisen mit 0,136 die mit Abstand höchste Streuung auf, gefolgt von den Schlüssel-Duplikaten mit 0,061 und der semantischen Länder-Vereinheitlichung mit 0,060. Die abgeleiteten Spalten und die einfache Bereinigung sind dagegen praktisch vollständig stabil. Der Zusammenhang ist plausibel: Wo eine Aufgabe eine eindeutige, formal beschreibbare Lösung besitzt, konvergieren wiederholte Generierungen auf dieselbe Umsetzung; wo hingegen eine Ermessensentscheidung wie die Wahl eines Ähnlichkeitsschwellenwerts zu treffen ist, schlägt sich die probabilistische Natur des Modells unmittelbar im Ergebnis nieder \autocite{zhaoSurveyLargeLanguage2026,hassaniRoleChatGPTData2023}.

Diese Kopplung hat eine unangenehme praktische Konsequenz, die über die reine Kennzahl hinausreicht: Instabilität und Fehleranfälligkeit treten gemeinsam auf. Gerade dort, wo ein Ergebnis besonders sorgfältig zu prüfen wäre, ist es zugleich am wenigsten verlässlich wiederholbar, sodass eine einmalige Abnahme eines Verarbeitungsschritts dessen künftiges Verhalten nicht garantiert. Für den produktiven Einsatz spricht dies dafür, den einmal erzeugten und geprüften Code als festes Artefakt zu übernehmen, statt ihn bei jeder Ausführung neu erzeugen zu lassen -- eine Konsequenz, die der Modus der Code-Generierung im Unterschied zur direkten Verarbeitung überhaupt erst zulässt. Die Unterschiede zwischen den Modellen sind demgegenüber nachrangig; die mittlere Streuung liegt zwischen 0,028 und 0,056, ohne dass sich daraus eine belastbare Rangfolge ergäbe.

\section{Bestätigungslauf auf einem unbekannten Datensatz}
\label{sec:bestaetigungslauf}

Alle bisher berichteten Ergebnisse beruhen auf dem Datensatz mit dem Startwert~1. An ihm wurden im Verlauf der Arbeit die Prompts formuliert, der Aufgabenzuschnitt festgelegt und die Auswertungslogik entwickelt; insbesondere die in Abschnitt~\ref{subsec:ergebnisse-bereinigung} beschriebene Angabe der Eingabe-Datentypen und die in Abschnitt~\ref{subsec:ergebnisse-duplikate} beschriebene Anweisung zum Umgang mit akademischen Graden sind unmittelbare Reaktionen auf dort beobachtete Fehler. Die an ihm gemessenen Werte sind damit nicht unbefangen: Sie beschreiben die Leistung auf einem Datensatz, gegen den die Konfiguration justiert wurde, und überschätzen die zu erwartende Leistung auf unbekannten Daten in unbekanntem Umfang. Diese Gefahr ist in der methodischen Literatur zur Bewertung von \acp{LLM} als zentrale Fehlerquelle beschrieben und wird dort regelmäßig als Grund dafür angeführt, dass berichtete Kennzahlen sich im Einsatz nicht bestätigen \autocite{boulahiaMulticriteriaEvaluationResearch2024,zhaoSurveyLargeLanguage2026}.

Um die Tragfähigkeit der Befunde zu prüfen, wurde deshalb ein zweiter, unabhängiger Datensatz mit dem Startwert~2 erzeugt und ohne jede weitere Anpassung von Prompts, Aufgabenstellungen oder Auswertungscode durchlaufen. Er folgt derselben Erzeugungsvorschrift aus Kapitel~\ref{ch:datensatz}, enthält jedoch andere Namen, andere Adressen und eine andere Zuordnung der eingebauten Abweichungen. Für den Direkt-Modus kommt die in Abschnitt~\ref{sec:ergebnisse-pipeline} eingeführte halbierte Fassung dieses Datensatzes zum Einsatz, die den Startwert~102 trägt, da die vollständige Fassung das Ausgabefenster der Modelle überschreitet. Weil die Erwartung an diesen Lauf vor seiner Ausführung feststand und nachträglich nichts mehr verändert wurde, trägt er die bestätigende Aussage dieser Arbeit, während der Startwert~1 die explorative Grundlage bildet.

\subsection{Verkettete und isolierte Ausführung}
\label{subsec:bestaetigung-pipeline}

Abbildung~\ref{fig:bestaetigung-pipeline} zeigt alle vier Kombinationen aus Ausführungsmodus und Eingabeform nebeneinander. Der ausgefüllte Balken gibt das Mittel über die Wiederholungen wieder, der hellere Aufsatz reicht bis zur besten Wiederholung; damit sind Durchschnitts- und Spitzenleistung gleichzeitig ablesbar, ohne dass der Mittelwert durch den Bestwert beschönigt würde. Schritte, die ein Lauf nicht mehr erreicht hat, gehen mit dem Wert null ein.

\begin{figure}[htbp]
    \centering
    \includegraphics[width=\textwidth]{bestaetigung_pipeline.png}
    \caption[Bestätigungslauf: Genauigkeit je Pipeline-Schritt]{Abgeglichene Genauigkeit je Schritt im Bestätigungslauf, getrennt nach Ausführungsmodus und Eingabeform. Balken: Mittel über die Wiederholungen; heller Aufsatz: beste Wiederholung.}
    \label{fig:bestaetigung-pipeline}
\end{figure}

Der Code-Generierungs-Modus bestätigt das Bild aus Abschnitt~\ref{sec:ergebnisse-pipeline} in allen wesentlichen Zügen. Über fünf Wiederholungen der verketteten Strecke erreichen die drei Cloud-Modelle in den Schritten~1 bis~8 durchweg Werte zwischen 96,0 und 100~Prozent und liegen damit im selben Band wie beim Startwert~1. Über alle neun Schritte gemittelt erzielt das Modell von Anthropic 96,1~Prozent, jenes von Google 94,9~Prozent und jenes von OpenAI 93,1~Prozent, gegenüber 94,2~Prozent für die regelbasierte Vergleichspipeline. Der Endschritt bricht erneut ein, und zwar auf 69,1, 60,6 und 49,1~Prozent bei einem Baseline-Wert von 53,0~Prozent. Auch die Rangfolge bleibt erhalten: Alle drei Modelle lösen die formal beschreibbaren Schritte praktisch fehlerfrei, und der zusammengesetzte Endschritt ist die einzige Stelle, an der die Kette substantiell verliert.

Der isolierte Durchlauf trennt auch hier eigene Leistung von geerbtem Fehler, und das Ergebnis fällt noch deutlicher aus als beim Startwert~1. Mit fehlerfreier Eingabe erreichen die Modelle über alle neun Schritte 99,7, 99,7 und 98,4~Prozent; den Endschritt lösen die Modelle von Anthropic und Google in jeder einzelnen Wiederholung vollständig korrekt, jenes von OpenAI erreicht konstant 90,9~Prozent. Der in der Kette gemessene Verlust von 30,9 bis 41,8~Prozentpunkten im Endschritt ist damit fast vollständig geerbt. Die zentrale Aussage aus Abschnitt~\ref{subsec:pipeline-fortpflanzung} -- dass Evaluierungen einzelner \ac{ETL}-Schritte und Evaluierungen ganzer Strecken systematisch in entgegengesetzte Richtungen verzerren -- ist damit an einem Datensatz repliziert, gegen den nichts justiert wurde \autocite{walhaDataIntegrationTraditional2024,boulahiaMulticriteriaEvaluationResearch2024}.

Für den Bestätigungslauf wurde zusätzlich die regelbasierte Pipeline isoliert ausgeführt, was in Abschnitt~\ref{sec:ergebnisse-pipeline} noch unterblieben war, und dieser Vergleich stellt den Befund auf eine breitere Grundlage. Ihre Werte sind allerdings asymmetrisch zu lesen: In den sieben formal beschreibbaren Schritten verwendet sie dieselben Routinen, aus denen auch die Kettenreferenz komponiert ist, sodass ein Ergebnis von 1,000 dort definitionsgemäß eintritt und keine Leistung ausweist. Aussagekräftig sind allein die beiden Schritte mit begrenzten Heuristiken, die sie mit 99,3 und 97,6~Prozent löst. Entscheidend ist jedoch der Endschritt: Auch die Regel löst ihn mit fehlerfreier Eingabe vollständig korrekt und fällt in der Kette auf 53,0~Prozent. Ihr Verlust von 47,0~Prozentpunkten übersteigt damit den jedes einzelnen Modells. Die Fehlerfortpflanzung ist folglich keine Eigenheit \ac{KI}-gestützter Verarbeitung, sondern eine Eigenschaft der verketteten Architektur, die das deterministische Verfahren in gleichem, hier sogar stärkerem Maße trifft \autocite{ridzuanReviewDataQuality2024,walhaDataIntegrationTraditional2024}.

Der Direkt-Modus liefert das gegenteilige Muster und bestätigt ebenfalls den Befund des Startwerts~101. In den Schritten~1 bis~8 arbeitet er auf dem halbierten Datensatz mit 96,3 bis 100~Prozent auf dem Niveau der Code-Generierung, wobei er die semantische Länder-Vereinheitlichung mit 98,8 bis 99,9~Prozent löst und damit die regelbasierte Pipeline mit 99,1~Prozent knapp erreicht oder übertrifft. Im Endschritt jedoch fällt er auf 2,7 bis 7,6~Prozent und damit weit unter jeden anderen Ansatz. Die Ursache ist keine inhaltliche: Die Antworten waren in keinem der Schritte abgeschnitten, alle neun Schritte lieferten eine vollständige, schemakonforme Tabelle. Entscheidend ist, dass der isolierte Durchlauf hier nichts rettet -- mit fehlerfreier Eingabe erreicht der Endschritt lediglich 1,5 bis 4,6~Prozent. Anders als bei der Code-Generierung ist dieser Einbruch also gerade nicht geerbt, sondern eigenes Versagen: Das Modell muss im Direkt-Modus die Aggregation über 60~Gruppen selbst rechnen, statt eine Rechenvorschrift zu formulieren, und genau diese arithmetische Zusammenfassung mehrerer hundert Einzelwerte gelingt ihm nicht. Der Modus vertauscht damit die Stärken: Er gewinnt dort, wo Einzelfallverständnis gefragt ist, und verliert dort, wo verlässliche Rechenoperationen über viele Datensätze verlangt werden \autocite{khanCognitiveETLUsing2025,hassaniRoleChatGPTData2023}.

\subsection{Reproduzierbarkeit über die Wiederholungen}
\label{subsec:bestaetigung-streuung}

Der Bestätigungslauf wurde im Code-Generierungs-Modus je fünfmal verkettet und isoliert ausgeführt, im Direkt-Modus je zweimal, wodurch sich die in Abschnitt~\ref{sec:ergebnisse-reproduzierbarkeit} auf Aufgabenebene betrachtete Reproduzierbarkeit erstmals auf Ebene ganzer Strecken beziffern lässt. Abbildung~\ref{fig:bestaetigung-streuung} zeigt jede einzelne Wiederholung als Punkt, links als Gesamtwert der Strecke und rechts für den Endschritt.

\begin{figure}[htbp]
    \centering
    \includegraphics[width=\textwidth]{bestaetigung_streuung.png}
    \caption[Reproduzierbarkeit im Bestätigungslauf]{Streuung über die Wiederholungen des Bestätigungslaufs. Jeder Punkt ist ein vollständiger Durchlauf, der waagerechte Strich das Mittel.}
    \label{fig:bestaetigung-streuung}
\end{figure}

Das Ergebnis ist der aufschlussreichste Einzelbefund des Bestätigungslaufs. Auf der Ebene des Gesamtwerts wirken die Läufe stabil: Die Standardabweichung über die fünf verketteten Wiederholungen beträgt lediglich 0,6~Prozentpunkte beim Modell von Anthropic, 2,0 bei Google und 3,2 bei OpenAI. Der Endschritt zeigt jedoch ein völlig anderes Bild. Dort schwankt das Ergebnis desselben Modells bei identischem Prompt und identischen Eingabedaten zwischen 27,3 und 72,7~Prozent bei OpenAI sowie zwischen 39,4 und 81,8~Prozent bei Google; lediglich das Modell von Anthropic bleibt mit 63,6 bis 72,7~Prozent in einem engeren Band. Eine einzelne Messung des Endergebnisses ist damit praktisch nicht aussagekräftig: Je nach Wiederholung ließe sich dasselbe Modell als klar besser oder klar schlechter als die regelbasierte Pipeline darstellen.

Entscheidend ist der Vergleich mit dem isolierten Durchlauf, denn dort verschwindet diese Streuung vollständig. Alle drei Modelle liefern über sämtliche fünf Wiederholungen hinweg exakt denselben Endschritt-Wert, und die Standardabweichung des Gesamtwerts sinkt auf 0,08 bis 0,62~Prozentpunkte -- bei den Modellen von Anthropic und Google um mehr als eine Größenordnung. Die beobachtete Instabilität entsteht folglich nicht aus der probabilistischen Natur der Textgenerierung selbst, sondern aus deren Verstärkung entlang der Kette: Eine kleine, für sich genommen belanglose Abweichung in der Wahl des Duplikat-Repräsentanten legt fest, welche Bestellungen im Endschritt einen Kundendatensatz finden, und schlägt dort um ein Vielfaches verstärkt durch. Die Reproduzierbarkeit einer \ac{KI}-gestützten \ac{ETL}-Strecke ist damit keine Eigenschaft des eingesetzten Modells, sondern eine Eigenschaft ihrer Architektur \autocite{ridzuanReviewDataQuality2024,zhaoSurveyLargeLanguage2026}.

Praktisch folgt daraus eine schärfere Fassung der bereits in Abschnitt~\ref{sec:ergebnisse-reproduzierbarkeit} formulierten Empfehlung. Es genügt nicht, den erzeugten Code als festes Artefakt zu übernehmen; ebenso festzuschreiben sind die Zwischenergebnisse der frühen Schritte, da bereits deren Variation ausreicht, um das Endergebnis um über 40~Prozentpunkte zu verschieben, ohne dass sich an einer der beteiligten Vorschriften etwas geändert hätte.

\subsection{Kosten, Token-Verbrauch und Laufzeit}
\label{subsec:bestaetigung-aufwand}

Abbildung~\ref{fig:bestaetigung-aufwand} stellt den Aufwand einer vollständigen Strecke in beiden Ausführungsmodi gegenüber. Die Werte sind Mittel über die Wiederholungen und beziehen sich jeweils auf einen kompletten Durchlauf aller neun Schritte.

\begin{figure}[htbp]
    \centering
    \includegraphics[width=\textwidth]{bestaetigung_aufwand.png}
    \caption[Aufwand der beiden Ausführungsmodi]{Kosten, Laufzeit und Token-Verbrauch einer vollständigen Strecke je Ausführungsmodus. Die Direktverarbeitung arbeitet dabei auf dem halbierten Datensatz.}
    \label{fig:bestaetigung-aufwand}
\end{figure}

Eine vollständige Strecke im Code-Generierungs-Modus kostet zwischen 0,176 und 0,216~US-Dollar und benötigt zwischen 102 und 135~Sekunden. Der Direkt-Modus liegt mit 1,32 bis 1,71~US-Dollar und 435 bis 854~Sekunden um den Faktor sieben beziehungsweise vier darüber -- und das, obwohl er nur den halbierten Datensatz verarbeitet. Auf die Zeile bezogen beträgt der Kostenunterschied damit rund das Fünfzehnfache. Der Grund liegt in der Struktur der Verarbeitung: Im Code-Generierungs-Modus durchläuft nur die Aufgabenbeschreibung das Modell, während die Daten lokal verarbeitet werden; im Direkt-Modus wandert die vollständige Tabelle in jedem der neun Schritte einmal hinein und einmal wieder heraus. Entsprechend steigt der Token-Umsatz von rund 15.000 bis 32.000 auf 169.000 bis 285.000 je Strecke.

Der Token-Verbrauch macht zugleich einen Unterschied zwischen den Modellen sichtbar, der sich in der Genauigkeit nicht niederschlägt. Das Modell von Google benötigt im Code-Generierungs-Modus mit rund 28.000~Ausgabe-Token annähernd das Dreifache der beiden anderen Modelle, obwohl seine erzeugten Skripte nicht länger sind. Die Ursache sind die verborgenen Zwischenschritte des Modells, die als Ausgabe abgerechnet, aber nicht ausgeliefert werden. Für die Kostenplanung ist dies erheblich, weil sich der Aufwand eines solchen Modells nicht aus der Länge der sichtbaren Antwort abschätzen lässt: Bei der Auslegung des Ausgabebudgets im Direkt-Modus führte genau diese Eigenschaft dazu, dass Antworten abbrachen, obwohl die erwartete Tabelle rechnerisch weit in das gewählte Budget gepasst hätte, und erforderte dessen Anhebung auf 96.000~Token \autocite{changEfficientPromptingMethods2024,2026AIIndex}.

In der Zusammenschau relativiert der Aufwand den in Abschnitt~\ref{subsec:staerken} beschriebenen Vorteil des Direkt-Modus erheblich. Sein Gewinn bei den semantischen Einzelschritten beträgt wenige Prozentpunkte, sein Verlust im zusammengesetzten Endschritt über 50 und sein Mehraufwand je verarbeiteter Zeile mehr als eine Größenordnung. Als durchgängiger Verarbeitungsmodus für eine vollständige \ac{ETL}-Strecke ist er damit nicht zu begründen; sinnvoll einsetzbar ist er allein als gezielt vorgeschalteter Einzelschritt für genau jene Teilaufgaben, die Einzelfallverständnis verlangen.

\subsection{Prompting-Strategien im Bestätigungslauf}
\label{subsec:bestaetigung-strategien}

Die dritte Untersuchungsebene des Bestätigungslaufs wiederholt die Faktormatrix aus Abschnitt~\ref{sec:vergleich-prompting} auf dem unbekannten Datensatz, beschränkt auf die drei klassischen Prompting-Strategien und die drei Cloud-Modelle. Die schema-angereicherte Variante wurde nicht erneut ausgeführt, da ihr in Abschnitt~\ref{sec:vergleich-prompting} belegter Einbruch nicht auf eine zufällige Schwankung, sondern auf einen benennbaren systematischen Mechanismus zurückgeht und eine Wiederholung folglich keinen zusätzlichen Erkenntnisgewinn verspricht. Es verbleiben 3~Strategien $\times$ 9~Aufgaben $\times$ 3~Modelle $\times$ 2~Wiederholungen, also 162~Läufe. Die Zahl der Wiederholungen wurde gegenüber dem explorativen Lauf von drei auf zwei verringert, da diese Ebene die Rangfolge der Strategien bestätigen soll und die Reproduzierbarkeit bereits in Abschnitt~\ref{sec:ergebnisse-reproduzierbarkeit} über 432~Läufe sowie in Abschnitt~\ref{subsec:bestaetigung-streuung} über fünf vollständige Streckendurchläufe bestimmt wurde; zwei Wiederholungen zeigen hier folglich an, ob ein beobachteter Strategieunterschied überhaupt Bestand hat, ohne die Streuungsaussage selbst tragen zu müssen.

Anders als bei der Pipeline ist hier kein gesonderter isolierter Durchlauf erforderlich, denn die Faktormatrix ist bereits ihrer Anlage nach isoliert: Jede Aufgabe liest unmittelbar die Rohdaten, und kein Ergebnis wird an eine folgende Aufgabe weitergereicht. Die in Abschnitt~\ref{subsec:pipeline-fortpflanzung} untersuchte Fehlerfortpflanzung kann in dieser Ebene folglich gar nicht auftreten, weshalb die gemessenen Unterschiede vollständig der jeweiligen Strategie zuzurechnen sind.

\begin{figure}[htbp]
    \centering
    \includegraphics[width=\textwidth]{bestaetigung_strategien.png}
    \caption[Prompting-Strategien im Bestätigungslauf]{Abgeglichene Genauigkeit der drei klassischen Prompting-Strategien je Aufgabe und Modell auf dem unbekannten Datensatz, gemittelt über zwei Wiederholungen. Die gestrichelte Linie markiert die regelbasierte Pipeline.}
    \label{fig:bestaetigung-strategien}
\end{figure}

Das Ergebnis bestätigt die in Abschnitt~\ref{sec:vergleich-prompting} vorgenommene Einordnung und widerlegt zugleich die dort zunächst berichtete Gesamtrangfolge. Über alle Aufgaben gemittelt führt auf dem unbekannten Datensatz das Few-Shot-Prompting mit 77,9~Prozent vor dem Chain-of-Thought-Prompting mit 75,1 und dem Zero-Shot-Prompting mit 71,8~Prozent -- also in genau umgekehrter Reihenfolge zu den 65,5, 64,1 und 60,9~Prozent des explorativen Laufs. Dieser Widerspruch ist jedoch kein Widerspruch der Sache, sondern eine Folge der Grundmenge: Der explorative Lauf umfasste das lokale Modell, der Bestätigungslauf aus den in Abschnitt~\ref{sec:vergleich-modelle} genannten Gründen nicht. Beschränkt man den explorativen Lauf ebenfalls auf die drei kommerziellen Modelle, liegen dessen Werte mit 72,6, 73,3 und 73,4~Prozent so eng beieinander, dass er zwischen den Strategien überhaupt nicht unterscheidet.

Was sich dagegen über beide Datensätze hinweg reproduziert, ist das aufgabenbezogene Muster, und zwar mit deutlich vergrößertem Abstand. Bei den unscharfen Duplikaten erreicht das Few-Shot-Prompting 82,0~Prozent gegenüber 40,0 für Chain-of-Thought und 26,6 für Zero-Shot; im explorativen Lauf lagen dieselben Werte bei 41,5, 21,3 und 30,7~Prozent. Bei den Schlüssel-Duplikaten führt erneut das Chain-of-Thought-Prompting mit 46,7 gegenüber jeweils 33,3~Prozent für die beiden anderen Strategien -- dieselbe Konstellation wie zuvor mit 51,1 gegenüber 33,3. Die übrigen sieben Aufgaben unterscheiden auch hier nicht zwischen den Strategien: Die Bereinigungsaufgaben und die exakten Duplikate werden von allen drei nahezu fehlerfrei gelöst, die zusammengesetzte Aufgabe aus Join, Bereinigung und Aggregation von keiner -- sie scheitert erneut in allen 18~Läufen. Die Wirkung der Prompt-Gestaltung ist damit reproduzierbar an die Aufgabenart gebunden, ihre Rangfolge über alle Aufgaben hinweg dagegen nicht \autocite{sahooSystematicSurveyPrompt,changEfficientPromptingMethods2024}.

\subsection{Prompting-Strategien ohne Fehlerfortpflanzung}
\label{subsec:bestaetigung-gruppen}

Die bis hierhin berichteten Strategieunterschiede unterliegen einer Einschränkung, die aus dem Aufbau der Faktormatrix folgt. Jede Aufgabe wird dort auf den unveränderten Rohdaten gestellt, sodass die anspruchsvollen Stufen zusätzlich alle vorgelagerten Mängel mitbehandeln müssen. Bei der zusammengesetzten Aufgabe führt das dazu, dass jede Strategie gleichermaßen auf null fällt -- eine Aufgabe, die alle Verfahren scheitern lässt, unterscheidet sie nicht mehr. Die Messung ist an dieser Stelle blind, und zwar nicht zufällig, sondern systematisch bei genau den Aufgaben, bei denen ein Strategieeffekt am ehesten zu erwarten wäre.

Um diese Blindstelle zu schließen, wurde eine vierte Untersuchungsebene ergänzt, die zwischen der Einzelaufgabe und der vollständigen Neunerkette liegt. Die neun Pipeline-Schritte werden zu drei Gruppen von je drei Schritten zusammengefasst -- Bereinigung, Deduplizierung und Transformation --, und jede Gruppe erhält die Soll-Ausgabe ihrer Vorgängergruppe als Eingabe. Die Anweisungen der drei zugehörigen Schritte werden dabei wörtlich aus der Pipeline übernommen und nummeriert in einem Prompt zusammengefasst, sodass die Gruppe inhaltlich dasselbe verlangt wie die drei Einzelschritte. Damit bleibt die Komposition mehrerer Verarbeitungsschritte als Anforderung erhalten -- genau jene Eigenschaft, an der sich eine Strategie wie das Chain-of-Thought-Prompting bewähren sollte --, ohne dass ein Fehler aus einer früheren Gruppe das Ergebnis überdeckt. Ausgeführt wurden 3~Strategien $\times$ 3~Modelle $\times$ 3~Gruppen, also 27~Läufe.

\begin{figure}[htbp]
    \centering
    \includegraphics[width=\textwidth]{bestaetigung_gruppen.png}
    \caption[Prompting-Strategien mit und ohne Fehlerfortpflanzung]{Dieselben Aufgaben als gedrittelte Strecke mit Soll-Eingabe je Gruppe (links) und als unabhängige Einzelaufgaben auf den Rohdaten (rechts), gemittelt über die drei Modelle.}
    \label{fig:bestaetigung-gruppen}
\end{figure}

Das Ergebnis ist eindeutig und in seiner Deutlichkeit unerwartet: Sämtliche 27~Läufe lieferten ein verwertbares Ergebnis, und die Genauigkeit liegt durchgängig zwischen 98,4 und 100~Prozent. Über die drei Gruppen gemittelt erreicht das Chain-of-Thought-Prompting 99,5~Prozent, das Zero-Shot-Prompting 99,3 und das Few-Shot-Prompting 99,0 -- eine Spanne von einem halben Prozentpunkt. Auch zwischen den Modellen ist kein belastbarer Unterschied mehr erkennbar. Besonders aufschlussreich ist die Transformationsgruppe: Sie enthält als letzten Schritt genau jene Aggregation, die als Einzelaufgabe in allen 18~Läufen scheiterte, und wird hier von jedem Modell in jeder Strategie und im ersten Versuch vollständig korrekt gelöst.

Abbildung~\ref{fig:bestaetigung-gruppen} stellt beide Betrachtungsweisen unmittelbar gegenüber und macht damit die eigentliche Konsequenz sichtbar. Der auffälligste Strategieunterschied der Faktormatrix -- 53, 72 und 62~Prozent bei der Duplikaterkennung -- verschwindet vollständig, sobald dieselben Schritte mit sauberer Eingabe und in einem klar abgegrenzten Zuschnitt gestellt werden. Was in der Matrix als Wirkung der Prompt-Gestaltung erscheint, ist demnach zu einem erheblichen Teil eine Wirkung des Aufgabenzuschnitts: Beispiele und schrittweise Herleitung helfen einem Modell, eine überladene Aufgabe zu bewältigen, sind aber entbehrlich, sobald die Aufgabe angemessen zugeschnitten ist. Für die Beantwortung von FF2 ist das ein Befund von eigenständigem Gewicht, denn er kehrt die übliche Reihenfolge der Gestaltungsentscheidungen um: Nicht der Prompt ist zuerst zu optimieren, sondern der Zuschnitt der Aufgabe -- und ist dieser einmal richtig gewählt, verliert die Wahl der Prompting-Strategie ihre Bedeutung weitgehend \autocite{khanOverviewETLTechniques2024,schulhoffPromptReportSystematic2025}.

Der Aufwand dieser Ebene unterstreicht den Befund zusätzlich. Die 27~Läufe kosteten zusammen 1,27~US-Dollar bei 13,1~Minuten Laufzeit, und von den 27~Läufen benötigten 25 nur einen einzigen Versuch. Dem stehen 162~Läufe der Faktormatrix mit 3,14~US-Dollar gegenüber, von denen 37 kein verwertbares Ergebnis lieferten. Der kleinteiligere Zuschnitt ist damit nicht nur genauer, sondern zugleich günstiger und zuverlässiger -- eine Kombination, die in dieser Untersuchung sonst an keiner Stelle auftritt.

\section{Vergleich KI-basiert vs. klassisch}
\label{sec:vergleich-baseline}

Abbildung~\ref{fig:uebersicht-modi} stellt die drei untersuchten Ansätze je Aufgabenkategorie gegenüber. Für die Code-Generierung ist dabei die jeweils beste klassische Prompting-Strategie ausgewiesen, sodass der Vergleich das jeweils Erreichbare und nicht den Durchschnitt einer beliebigen Konfiguration abbildet.

\begin{figure}[htbp]
    \centering
    \includegraphics[width=0.9\textwidth]{uebersicht_modi.png}
    \caption[Ansätze im Vergleich je Kategorie]{Mittlere Genauigkeit der drei Ansätze je Aufgabenkategorie. Für die Code-Generierung ist jeweils die beste klassische Prompting-Strategie ausgewiesen.}
    \label{fig:uebersicht-modi}
\end{figure}

Das Ergebnis fällt für die \ac{KI}-gestützten Verfahren zunächst ernüchternd aus: Die regelbasierte Pipeline liegt in allen drei Kategorien vorn, bei der Datenbereinigung mit 98,5 gegenüber 89,3 und 56,9~Prozent, bei der Duplikaterkennung mit 99,0 gegenüber 69,0 und 66,2~Prozent und bei der Datentransformation mit 70,7 gegenüber 60,1 und 54,1~Prozent. Dieser Befund ist jedoch differenziert zu lesen. Für die sechs formal beschreibbaren Aufgaben ist er weitgehend definitorisch bedingt: Wie in Abschnitt~\ref{sec:klassische-pipeline} dargelegt, verwendet die regelbasierte Pipeline für diese Aufgaben dieselben deterministischen Routinen, die auch der Erzeugung der Soll-Lösung zugrunde liegen; die explizit formulierte Regel \emph{ist} dort die korrekte Lösung. Ein Ergebnis von 100~Prozent ist folglich kein Leistungsnachweis, sondern die zu erwartende Untergrenze für ein Verfahren, dessen Vorschrift mit der Aufgabendefinition zusammenfällt. Interessant ist an diesen Aufgaben nicht, dass die Regel gewinnt, sondern dass die Modelle mit gut gewähltem Prompt gleichauf liegen, ohne dass die Regel zuvor formuliert werden musste.

Aussagekräftig ist der Vergleich daher erst bei den drei Aufgaben, für die die regelbasierte Pipeline bewusst mit begrenzten Heuristiken ausgestattet wurde. Bei der Länder-Vereinheitlichung erreicht sie 95,5~Prozent: Ihre Zuordnungstabelle deckt ISO-Codes in zwei- und dreistelliger Form, IOC-Kürzel sowie deutsche und englische Ländernamen ab, scheitert aber an 42 von 540~Datensätzen, nämlich an \texttt{England} (17), \texttt{Holland} (14) und dem Tippfehler \texttt{Deutshcland} (11). Genau hier liegt die Trennlinie zwischen beiden Verfahrensarten: Eine Zuordnungstabelle kann nur erfassen, was zuvor jemand aufgezählt hat, während ein Sprachmodell erschließt, dass \enquote{Holland} umgangssprachlich für die Niederlande steht. Dass die Code-Generierung diesen Vorteil dennoch nicht ausspielt, liegt daran, dass sie das Weltwissen des Modells in genau dieselbe aufzählende Form zwingt und dabei ebenfalls Lücken lässt -- lediglich andere. Erst der Direkt-Modus, in dem jeder Wert einzeln beurteilt wird, übertrifft die regelbasierte Lösung mit 96,4 bis 96,5~Prozent.

Noch klarer zeigt sich dieses Muster bei den unscharfen Duplikaten. Die regelbasierte Pipeline erreicht hier 96,9~Prozent, weil ihr normalisierter Schlüssel aus Name und E-Mail-Bestandteil genau jene Abweichungen erfasst, die durch Groß- und Kleinschreibung, zusätzliche Leerzeichen oder Punkte im lokalen Teil der Adresse entstehen -- und der Datensatz eben diese Muster enthält. Die Code-Generierung bleibt mit 69,0~Prozent deutlich zurück, der Direkt-Modus erreicht mit bis zu 96,5~Prozent nahezu das Niveau der Regel. Diese Konstellation macht zugleich eine Grenze der Untersuchungsanlage sichtbar, die in Kapitel~\ref{ch:fazit} aufzugreifen ist: Die eingebauten Abweichungen entstammen einer bekannten, endlichen Menge von Mustern, was regelbasierte Verfahren systematisch begünstigt. In realen Beständen, deren Abweichungsmuster nicht vorab bekannt sind, verschiebt sich dieses Verhältnis zugunsten der semantischen Verfahren \autocite{coteDataCleaningMachine2024,ridzuanReviewDataQuality2024}.

Der dritte aussagekräftige Fall ist die zusammengesetzte Aufgabe aus Join, Bereinigung und Aggregation, bei der auch die regelbasierte Pipeline mit 12,1~Prozent nahezu vollständig scheitert -- und zwar aus derselben Ursache wie die Modelle: Ihre begrenzte Zuordnungstabelle lässt 57 Länderangaben als unbekannt zurück, die anschließend eine eigene, in der Soll-Lösung nicht vorgesehene Gruppe bilden und die Aggregation über alle betroffenen Kategorien verfälschen. Die Aufgabe trennt somit nicht zwischen den Ansätzen, sondern zeigt, dass ein Fehler in einem frühen semantischen Schritt in einer zusammengesetzten Verarbeitung unabhängig vom gewählten Verfahren dieselbe verstärkende Wirkung entfaltet. In der zerlegten Kette hingegen lag die regelbasierte Pipeline mit 60,6~Prozent im Endschritt hinter allen drei Cloud-Modellen, weil sie mehr Bestellungen ohne zugehörigen Kunden zurückließ.

Zusammengefasst ergibt sich kein Verfahren, das über alle Aufgaben hinweg vorzuziehen wäre. Für formal beschreibbare Aufgaben ist die regelbasierte Umsetzung der \ac{KI}-gestützten mindestens ebenbürtig und in Kosten, Geschwindigkeit und Determiniertheit deutlich überlegen -- 0,17~Sekunden ohne Nutzungskosten gegenüber gut zwei Minuten und rund 0,19~US-Dollar für dieselbe Strecke. Ihr Vorteil endet jedoch genau dort, wo die Verarbeitungsvorschrift nicht mehr vollständig vorab formulierbar ist, und in diesem Bereich liegt der Beitrag der Modelle weniger im Ersetzen als im Ergänzen der Regel \autocite{khanLeveragingArtificialIntelligence2025,kasireddyRoleAIModern2026}.

\section{Diskussion der Ergebnisse}
\label{sec:diskussion}

\subsection{Stärken des LLM-Einsatzes}
\label{subsec:staerken}

Die deutlichste Stärke der untersuchten Modelle liegt nicht in der erreichten Genauigkeit, sondern im Aufwand, mit dem sie erreicht wird. Die regelbasierte Vergleichspipeline setzt voraus, dass sämtliche Verarbeitungsvorschriften einschließlich der Zuordnungstabelle für Länderangaben und der Normalisierungsregel für Duplikate zuvor formuliert werden; ihre hohe Ergebnisqualität ist das Resultat dieser vorgelagerten Arbeit. Die Modelle erreichen bei den formal beschreibbaren Aufgaben mit gut gewähltem Prompt vergleichbare Werte, ohne dass eine solche Spezifikation vorliegt, und zwar in Sekunden und zu Kosten im Bereich weniger Cent. Für Vorhaben, in denen die Verarbeitungslogik noch nicht feststeht oder sich häufig ändert, verschiebt dies die Aufwandsrechnung erheblich \autocite{khanLeveragingArtificialIntelligence2025,aAIEnhancedETLProcesses2025}.

Die zweite Stärke betrifft die semantisch geprägten Aufgaben, allerdings in einer präziseren Form als zu Beginn erwartet. Der Vorteil des Sprachverständnisses tritt nur dann zutage, wenn das Modell sein Wissen auf den einzelnen Wert anwenden kann, wie es der Direkt-Modus zulässt: Bei den unscharfen Duplikaten steigt die Genauigkeit dort von 69,0 auf bis zu 96,5~Prozent, bei der Länder-Vereinheitlichung von 95,8 auf 96,5~Prozent und damit erstmals über die regelbasierte Lösung. Sobald dieselbe Aufgabe in generierten Code gefasst werden muss, geht dieser Vorteil weitgehend verloren, weil das Modell sein Wissen in eine aufzählende Tabelle oder einen numerischen Schwellenwert übersetzen muss und damit auf die Ausdrucksmittel des regelbasierten Verfahrens zurückfällt. Der Ausführungsmodus ist somit kein technisches Detail, sondern bestimmt maßgeblich, ob die semantische Fähigkeit des Modells überhaupt wirksam werden kann \autocite{hassaniRoleChatGPTData2023,khanCognitiveETLUsing2025}.

Als dritte Stärke erweist sich die Zerlegbarkeit der Aufgabenstellung. Der Vergleich zwischen verketteter Pipeline und unabhängigem Einzellauf in Abbildung~\ref{fig:pipeline-vs-einzellauf} zeigt, dass dieselben Modelle bei kleinteilig gestellten Schritten erheblich besser abschneiden, in einzelnen Schritten um über 50~Prozentpunkte. Die vollständig gescheiterte zusammengesetzte Aufgabe wird in zerlegter Form von allen vier Modellen fehlerfrei gelöst. Nicht die Schwierigkeit der einzelnen Operation begrenzt die Leistungsfähigkeit, sondern die Zahl der Entscheidungen, die in einem Arbeitsschritt gleichzeitig zu treffen sind -- eine Eigenschaft, die sich durch den Zuschnitt der Aufgabenstellung unmittelbar beeinflussen lässt.

\subsection{Schwächen und Risiken}
\label{subsec:schwaechen}

Das gravierendste Risiko ist das stille Versagen. Der in Skript~\ref{lst:pandas-object} dokumentierte Fall zeigt einen erzeugten Code, der fehlerfrei durchläuft, eine schemakonforme Ausgabe erzeugt und dennoch nichts bewirkt. Weder die Erfolgsquote noch eine Schemaprüfung noch eine Sichtprüfung des Codes hätten diesen Fehler aufgedeckt; sichtbar wurde er allein durch den Abgleich mit einer bekannten Soll-Lösung, die im produktiven Einsatz gerade nicht vorliegt. Verschärfend kommt hinzu, dass die Ursache -- der Rückgriff auf ein in den Trainingsdaten dominantes, inzwischen überholtes Idiom -- systematisch und damit reproduzierbar ist. Jede Produktivnutzung generierten \ac{ETL}-Codes setzt folglich eine unabhängige Ergebnisprüfung voraus, die nicht auf der Abwesenheit von Fehlermeldungen beruhen darf \autocite{hassaniRoleChatGPTData2023,mumuniAutomatedDataProcessing2025}.

Ein zweites Risiko liegt in der Fehlerfortpflanzung. Der isolierte Kettendurchlauf hat gezeigt, dass ein Schritt, der für sich betrachtet 94,6~Prozent erreicht und damit unauffällig wirkt, drei Schritte später einen vollständigen Abbruch auslösen kann. Die üblichen Qualitätskennzahlen eines einzelnen Schritts sind daher kein verlässlicher Indikator für die Verwendbarkeit seines Ergebnisses in nachgelagerten Schritten, da nicht der Umfang, sondern die Art der verbliebenen Abweichung entscheidet. Ein Ergebnis mit wenigen, aber typfremden Restwerten kann schädlicher sein als eines mit vielen, aber gleichartigen Abweichungen.

Das dritte Risiko ist die begrenzte Reproduzierbarkeit, die sich, wie in Abschnitt~\ref{sec:ergebnisse-reproduzierbarkeit} gezeigt, ausgerechnet bei den schwierigsten Aufgaben am stärksten ausprägt. Hinzu tritt die Abhängigkeit von serverseitig gepflegten Modellen, deren Verhalten sich ohne Zutun der Anwendenden ändern kann; diesem Umstand wurde durch die Protokollierung der exakten Versionskennung Rechnung getragen, er lässt sich damit jedoch nur dokumentieren und nicht beheben \autocite{zhaoSurveyLargeLanguage2026,2026AIIndex}.

Schließlich ist die Erwartung zu korrigieren, mehr Kontext im Prompt führe zu besseren Ergebnissen. Die schema-angereicherte Strategie hat die Genauigkeit bei der Duplikaterkennung auf ein Viertel der übrigen Strategien gesenkt, und das aufwändige Few-Shot-Prompting verursachte den dreifachen Token-Aufwand ohne Qualitätsgewinn. Die Gestaltung des Prompts ist damit ein zweischneidiges Werkzeug, dessen Wirkung sich nicht aus dem investierten Aufwand ableiten lässt, sondern empirisch je Aufgabenart zu bestimmen ist \autocite{schulhoffPromptReportSystematic2025,changEfficientPromptingMethods2024}.

\subsection{Einordnung in den Stand der Forschung}
\label{subsec:einordnung}

Die Ergebnisse bestätigen die in Kapitel~\ref{ch:stand-der-forschung} referierte grundsätzliche Einschätzung, dass \acp{LLM} für Aufgaben der Datenbereinigung und -integration nutzbar sind, präzisieren sie jedoch in einem für die Praxis wesentlichen Punkt. Während die Literatur den Nutzen überwiegend konzeptionell oder anhand einzelner Anwendungsfälle beschreibt, zeigt der hier durchgeführte kontrollierte Vergleich, dass die Eignung nicht vom Modell, sondern vom Zusammenspiel aus Aufgabentyp, Ausführungsmodus und Aufgabenzuschnitt abhängt \autocite{khanLeveragingArtificialIntelligence2025,hassaniRoleChatGPTData2023}. Die Spreizung zwischen den Aufgaben übersteigt jene zwischen den Modellen um ein Vielfaches, was der verbreiteten Orientierung an Modellranglisten eine begrenzte Aussagekraft für konkrete \ac{ETL}-Vorhaben bescheinigt \autocite{2026AIIndex,ComparisonAIModels}.

Hypothese~H1 gilt damit in ihrem Kern als bestätigt und in einem Teilaspekt als eingeschränkt bestätigt. Bestätigt ist die starke Abhängigkeit vom Aufgabentyp, die Ebenbürtigkeit bei formal beschreibbaren Aufgaben, der Verlust an Zuverlässigkeit mit zunehmender struktureller Komplexität sowie die Aussage, dass kein Ansatz über alle Aufgaben hinweg überlegen ist. Der erwartete Vorteil bei semantisch geprägten Aufgaben trat hingegen nur im Direkt-Modus ein, während er im primären Modus der Code-Generierung ausblieb -- eine Präzisierung, die in der Literatur bislang kaum Beachtung findet, weil dort die Wahl des Ausführungsmodus selten als eigenständige Gestaltungsdimension behandelt wird \autocite{khanCognitiveETLUsing2025,nagabhyruBridgingTraditionalETL2022}.

Hypothese~H2 ist in ihrer allgemeinen Form zu verwerfen. Aufwändigere Prompting-Strategien verbessern die Ergebnisqualität gegenüber dem Zero-Shot-Prompt nicht durchgängig; über alle Läufe hinweg schneidet der einfachste Prompt am besten ab, und die aufwändigste Variante fällt deutlich zurück. Zutreffend bleibt der erwartete abnehmende Grenznutzen, allerdings in verschärfter Form: Der Nutzen ist nicht nur abnehmend, sondern bei unpassendem Zuschnitt negativ. Bestätigt hat sich zugleich der in der Literatur beschriebene aufgabenabhängige Charakter der Wirkung, denn das Chain-of-Thought-Prompting führte bei den mehrstufigen Aufgaben der Duplikaterkennung und Transformation, während es bei den einstufigen Bereinigungsaufgaben schadete \autocite{sahooSystematicSurveyPrompt,liuPretrainPromptPredict2023}.

Über die beiden Hypothesen hinaus liefert die verkettete Betrachtung einen Beitrag, der in der Literatur bislang wenig ausgeprägt ist. Die getrennte Messung von eigener Schrittleistung und geerbtem Fehler zeigt, dass ein erheblicher Teil der in Einzelbetrachtungen berichteten Schwächen gar nicht dem bewerteten Schritt zuzurechnen ist: Der Endschritt der Kette, der in der verketteten Ausführung zwischen 71,2 und 81,8~Prozent erreicht, ist mit fehlerfreier Eingabe von allen vier Modellen vollständig lösbar. Evaluierungen, die einzelne \ac{ETL}-Operationen isoliert betrachten, überschätzen die Leistungsfähigkeit in Bezug auf reale Strecken daher ebenso, wie rein verkettete Betrachtungen sie unterschätzen; erst die Gegenüberstellung beider Betrachtungsweisen erlaubt eine belastbare Zuordnung \autocite{boulahiaMulticriteriaEvaluationResearch2024,walhaDataIntegrationTraditional2024}.

Ein methodischer Beitrag ergibt sich schließlich aus dem in Abschnitt~\ref{sec:bestaetigungslauf} berichteten Bestätigungslauf. Er zeigt, dass die inhaltlichen Befunde auf einem unbekannten Datensatz Bestand haben, macht zugleich aber sichtbar, wie wenig belastbar Einzelmessungen an verketteten Strecken sind: Der Endschritt derselben Konfiguration schwankte über fünf identische Wiederholungen um bis zu 45~Prozentpunkte, während er im isolierten Durchlauf in jeder Wiederholung exakt denselben Wert lieferte. Vergleichende Bewertungen von \acp{LLM} in \ac{ETL}-Strecken, die auf einem einzelnen Durchlauf beruhen -- wie es in der referierten Literatur überwiegend der Fall ist --, messen unter diesen Bedingungen zu einem erheblichen Teil Zufall statt Leistungsfähigkeit \autocite{boulahiaMulticriteriaEvaluationResearch2024,2026AIIndex}.

Auf dieser Grundlage lassen sich im folgenden Kapitel~\ref{ch:handlungsempfehlungen} konkrete Empfehlungen für die Auswahl zwischen klassischem, \ac{KI}-gestütztem und hybridem Vorgehen je Aufgabentyp ableiten, wie sie FF1 einfordert, ebenso wie Empfehlungen zur Gestaltung der Prompts im Sinne von FF2.
